% =====================================================================
%  2 Desenvolvimento -- visão geral do ambiente
% =====================================================================
\section{Desenvolvimento}\label{sec:desenvolvimento}

As práticas foram feitas em um computador com Windows, usando o \textbf{Vagrant} para criar as
máquinas virtuais e o \textbf{VMware Workstation} para executá-las. Todas as VMs usam o Rocky Linux 9
(\textit{box} \texttt{generic/rocky9}) e estão ligadas na rede privada \texttt{192.168.1.0/24}. A
Tabela~\ref{tab:maquinas} resume as máquinas criadas ao longo da disciplina. Os arquivos completos
do laboratório estão versionados no repositório do trabalho, na pasta \texttt{relatorio/codigos}
(\url{\repositorio}); senhas e chaves foram substituídas por marcadores.

\begin{table}[H]
  \centering
  \caption{Máquinas virtuais criadas no laboratório}\label{tab:maquinas}
  \footnotesize
  \setlength{\tabcolsep}{5pt}
  \renewcommand{\arraystretch}{1.3}
  \begin{tabularx}{\linewidth}{>{\ttfamily}l >{\ttfamily}l l c >{\raggedright\arraybackslash}X}
    \toprule
    \normalfont\bfseries Hostname & \normalfont\bfseries IP privado &
    \bfseries Porta redirecionada (VM $\rightarrow$ host) & \bfseries RAM/vCPU & \bfseries Função \\
    \midrule
    \normalfont(padrão) & \normalfont-- (só NAT) & 80 $\rightarrow$ 127.0.0.1:8080 & 1 GB / 2 & Laboratório 1: Apache servindo uma página HTML \\
    control-node & 192.168.1.2 & -- & 1 GB / 2 & \textit{Control Node} do Ansible \\
    app01 & 192.168.1.3 & 80 $\rightarrow$ 8080 & 1 GB / 2 & Aplicação Java \textit{Easy Notes} (porta 8080) \\
    db01 & 192.168.1.4 & 3306 $\rightarrow$ 3306 & 1 GB / 2 & Banco de dados MariaDB \\
    docker-notes & 192.168.1.5 & 8080 $\rightarrow$ 127.0.0.1:8082 & 2 GB / 2 & Docker: aplicação em container \\
    sonarqube & 192.168.1.6 & 9000 $\rightarrow$ 127.0.0.1:9000 & 4 GB / 2 & SonarQube e SonarScanner \\
    jenkins & 192.168.1.7 & 8080 $\rightarrow$ 127.0.0.1:8081 & 4 GB / 2 & Jenkins (CI/CD) \\
    \bottomrule
  \end{tabularx}
  \fonte{Elaborado pelo autor (2026).}
\end{table}
